\subsection*{What this proof of concept shows}
\begin{itemize}
\item \textbf{The problem is well posed and fast to simulate.} A daily
  OHQ reservoir model, driven by observed unimpaired inflow, runs at
  about 1400 simulated days per second on 24 cores, with mass balance
  closed to $10^{-7}$. Functional flows can be scored automatically from
  CEFF-style targets derived from the same river's natural hydrology.
\item \textbf{Prescribed functional flows work, at a drought cost.} A
  static functional-flow hydrograph scores well (0.93), but it spends
  storage early in a drought without knowing how long the drought will
  last. This is the opening an adaptive operation should exploit.
\item \textbf{Off-the-shelf PPO does not exploit it here.} Even started
  exactly from the prescriptive rule, PPO drifted toward over-releasing.
  The extra release protects a noisy exploring policy, and its supply
  cost arrives months later. The improvement an adaptive rule could make
  is small compared with year-to-year hydrologic noise, so policy-gradient
  learning with per-day Gaussian exploration has a poor signal to work
  with.
\end{itemize}

\subsection*{How to get an adaptive rule that beats the prescription}
The results point to methods that evaluate a policy \emph{without}
exploration noise, on whole multi-year windows:
\begin{enumerate}
\item \textbf{Direct policy search.} Parameterize an adaptive rule with
  10--50 parameters: e.g.\ pulse magnitudes and baseflow levels as
  functions of storage and the runoff forecast, or a small network. Tune
  the parameters with an evolution strategy (CMA-ES) on the deterministic
  multi-year objective. This is the evolutionary multi-objective direct
  policy search (EMODPS) used in reservoir operations. It needs no
  per-day credit assignment and gives Pareto fronts directly.
\item \textbf{Value-based RL with discrete releases} (the outline's DQN).
  Its exploration (occasional random actions) does not bias the policy in
  the same way. It is worth trying with a few release levels \emph{relative
  to} the prescribed schedule.
\item \textbf{A larger signal.} The headroom grows with harder conditions:
  stricter targets, higher demand, or climate-perturbed inflows with
  earlier snowmelt. Under those conditions a static schedule fails more
  often, and adaptivity is worth more.
\end{enumerate}
